Portfolio optimization is a central problem in quantitative finance. Classical methods such as the Markowitz approach provide widely accepted solutions, while more recent alternatives such as Hierarchical Risk Parity (HRP) introduce approaches based on clustering and correlation distances.

Quantum computing opens new possibilities for the representation and processing of financial information. In this context, this work proposes the implementation and experimental validation of the Quantum Hierarchical Risk Parity (QHRP) algorithm, which integrates quantum data encoding through a variational ansatz with hierarchical risk ordering techniques. The objective is to assess whether this formulation enables alternative representations of the correlation space and to compare its behaviour with the classical HRP method.

The formulation of the quantum ansatz, the angular encoding scheme for financial data, the circuit architecture, and the hybrid processing workflow are presented. Experiments are carried out using quantum simulators and compared with classical implementations, evaluating the stability of the encoding, the reproducibility of the results, and the limitations of the approach.

\keywords{Quantum computing, portfolio optimization, Hierarchical Risk Parity, quantum encoding, financial correlation, experimental validation}